\documentclass{beamer}
\usefonttheme{professionalfonts}
%\usetheme{Rochester}
%\usetheme{AnnArbor}
\usetheme{Szeged}
\useoutertheme{infolines}
\usecolortheme{dolphin}
\usepackage{graphicx}
% \usepackage[english]{babel}
\usepackage{hyperref}
\usepackage{multicol}
\usepackage{color}
%\usepackage{smartdiagram}


%\definecolor{carmine}{rgb}{0.59, 0.0, 0.09}
\usepackage{tikz}
\usetikzlibrary{shapes.callouts}
\usepackage{multirow}

\usepackage[utf8]{inputenc}

\usepackage{enumerate}
\usepackage[shortlabels]{enumitem}
\usepackage{amsmath,amsthm,amssymb,amsfonts,mathrsfs,latexsym,stmaryrd}
\usepackage{tcolorbox}

\usepackage{color} 
\usepackage{amssymb, amsmath, amsbsy}
\usepackage{adjustbox,lipsum}
\usepackage{placeins}

\usepackage{textpos}
\usepackage{wrapfig}
\usepackage[spanish, es-tabla]{babel}

\usepackage{caption}
\setlength{\parindent}{15pt}
% Beamer already loads hyperref; avoid reloading it with conflicting options.

\usepackage{url}
\usepackage{float}
\usepackage{tikz}
\usetikzlibrary{patterns}


\title[ICS-294]{Econometría ICS-294\\ \small{Clase 1}}

\author{Francisco Alfaro Medina}
\institute[USM]
{Universidad Técnica Federico Santa María\\ % Your institution for the title page
	\medskip
	\textit{} % Your email address
	
}
\logo{

\includegraphics[width=1cm]{../Imagenes/Logo_UTFSM.png}
}
\date{}
\newcommand\Fontvi{\fontsize{9}{7.2}\selectfont}
\begin{document}
	
\begin{frame}
	\titlepage % Print the title page as the first slide

\end{frame}
\section{Introducción}

\begin{frame}{Objetivos de la clase}
\begin{itemize}
    \item Entender qué estudia la econometría y para qué se usa.
    \item Distinguir entre correlación, causalidad y predicción.
    \item Reconocer los principales tipos de datos.
    \item Repasar conceptos básicos de muestreo e inferencia.
\end{itemize}
\end{frame}

\begin{frame}{¿Qué es econometría?}
\begin{block}{Idea central}
La econometría usa teoría económica, datos y métodos estadísticos para responder preguntas económicas.
\end{block}

\begin{itemize}
    \item Contrasta teorías con evidencia.
    \item Permite pronosticar variables y evaluar políticas públicas.
    \item Traduce preguntas económicas en modelos que podemos estimar.
\end{itemize}
 
\end{frame}

\begin{frame}{Preguntas que queremos responder}
\begin{itemize}
    \item ¿Cuál es la brecha salarial por género?
    \item ¿Cuál es el efecto de completar la universidad sobre el salario?
    \item ¿Qué variables ayudan a predecir la inflación, el empleo o las ventas?
\end{itemize}

\begin{block}{Clave del curso}
No basta con encontrar una relación en los datos: debemos preguntarnos qué significa y bajo qué supuestos podemos interpretarla.
\end{block}
\end{frame}


\begin{frame}{Econometría moderna}


  La econometr\'ia moderna se enfoca en contestar 3 tipos de preguntas:
  \begin{itemize}
      \item Estimar como la esperanza condicional de una variable depende de otra.
      \item Estimar el impacto causal de una variable sobre la otra.
      \item Predecir el comportamento de una variable hacia el futuro.
  \end{itemize}

\end{frame}

\begin{frame}{M\'etodo econom\'etrico}
    
\begin{itemize}
    \item Definir una pregunta: ¿qué quiero estudiar?
    \item Plantear un modelo económico:
    \[
        Y = f(x_1,x_2,\ldots,x_k)
    \]
    \item Llevarlo a un modelo econométrico:
    \[
        y = f(x_1,x_2,\ldots,x_k) + u
    \]
    \item Usar datos para estimar, contrastar e interpretar.
\end{itemize}

\end{frame}

\begin{frame}{Ejemplo: retornos de la educación}
\begin{itemize}
    \item Pregunta: ¿cómo se relaciona la educación con el salario?
    \item Modelo económico:
    \[
        salario = f(educ, exp, ingfam, habilidad)
    \]
    \item Modelo econométrico:
    \[
        salario_i = \beta_0 + \beta_1 educ_i + u_i
    \]
\end{itemize}

\begin{block}{Cuidado}
El término $u_i$ resume factores no observados. Si esos factores están relacionados con educación, la interpretación causal se vuelve difícil.
\end{block}
\end{frame}

\begin{frame}{Causalidad y correlación}

{\color{blue}    Correlaci\'on no es causalidad}
\begin{itemize}
    \item  Muchas veces nos interesa conocer el efecto causal de una variable sobre otra.
    \item  El concepto {\color{blue} ceteris paribus } es clave para la causalidad: efecto de X sobre Y
{\it manteniendo todos los demas factores relevantes constantes}.
 Muy dif\'icil en la pr\'actica en ciencias sociales.
 \item  Vamos a ver en el curso los supuestos necesarios y t\'ecnicas para
hacerlo.
\end{itemize}
{\color{blue}Correlaci\'on espuria:} existe una relación matem\'atica pero los factores no
presentan una relaci\'on l\'ogica.
\end{frame}

\begin{frame}{Ejemplo: universidad e ingresos}
\begin{itemize}
\item Información: en una base de individuos de 35 años en Chile, quienes terminaron la universidad ganan 30\% más que quienes no la terminaron.
\item Pregunta: ¿podemos concluir que ir a la universidad causa un aumento del ingreso?
\end{itemize}

\begin{block}{Para discutir}
¿Qué otros factores podrían explicar esa diferencia?
\end{block}
\end{frame}


\begin{frame}{Ejemplo: interpretación}

No necesariamente podemos concluir causalidad.
\begin{itemize}
    \item Ir a la universidad está asociado con otras características que también afectan ingresos.
    \item Ejemplos: habilidad, redes, ingreso familiar, colegio de origen, comuna, motivación.
    \item La econometría busca separar estas explicaciones usando datos y supuestos explícitos.
\end{itemize}
    
\end{frame}


\begin{frame}{Econometr\'ia como herramienta aplicada}

\begin{itemize}
    \item {\color{blue}Salud y educación:} Miguel y Kremer (2004) estudian programas de desparasitación en colegios y encuentran mejoras en salud y asistencia escolar.
    \item {\color{blue}Productividad:} Bloom et al. (2017) analizan si la consultoría de gestión mejora la productividad de firmas en India.
    \item En ambos casos, la pregunta no es solo si hay una relación, sino cuánto cambia el resultado cuando cambia una política o intervención.
\end{itemize}
    \end{frame}


\begin{frame}{Tipos de datos: origen de la información}
      

\begin{itemize}
    \item {\color{blue}Experimentales}: el investigador asigna aleatoriamente la exposición o tratamiento.
    \item {\color{blue}Pseudo-experimentales}: hay una intervención o regla de asignación, pero no necesariamente aleatoria.
    \item {\color{blue}Observacionales}: el investigador observa lo que ocurre; no controla la exposición.
\end{itemize}

\begin{block}{Comentario}
En economía muchas veces trabajamos con datos observacionales, porque los experimentos pueden ser costosos, difíciles o éticamente problemáticos.
\end{block}
\end{frame}

\begin{frame}{Tipos de datos: corte transversal}

\begin{itemize}
    \item Observaciones de muchas unidades en un momento del tiempo.
    \item Ejemplos: individuos, hogares, empresas, regiones o países.
    \item Pregunta típica: ¿qué características explican diferencias entre personas o empresas?
\end{itemize}

\begin{block}{Ejemplo}
Encuesta CASEN: ingreso, educación, empleo y características del hogar en un año determinado.
\end{block}
\end{frame}


\begin{frame}{Tipos de datos: series de tiempo}

\begin{itemize}
    \item Observaciones de una o más variables a través del tiempo.
    \item Ejemplos: IPC, dólar, PIB, ventas mensuales, precio de acciones.
    \item El pasado puede influir en el presente: rezagos, ciclos y tendencias importan.
    \item La frecuencia puede ser diaria, mensual, trimestral, anual, etc.
\end{itemize}

\end{frame}

\begin{frame}{Tipos de datos: panel}

\begin{itemize}
    \item Seguimos a las mismas unidades durante varios períodos.
    \item Ejemplo: salarios, educación y empleo de las mismas personas durante 10 años.
    \item Combina variación entre unidades y variación en el tiempo.
\end{itemize}

\begin{block}{Ventaja}
Ayuda a controlar características no observadas que son relativamente estables en el tiempo.
\end{block}
\end{frame}

\section{Conceptos } 

\begin{frame}{Muestreo: población, unidad y muestra}

\begin{itemize}
    \item {\bf Unidad de observación:} individuo, hogar, empresa, comuna u objeto que se mide.
    \item {\bf Población:} conjunto total sobre el cual queremos aprender.
    \item {\bf Muestra:} subconjunto de la población que efectivamente observamos.
\end{itemize}

\begin{block}{Objetivo}
Usar la muestra para aprender sobre la población sin observarla completa.
\end{block}
 
\end{frame}




 \begin{frame}{Muestreo}
     Se dice que una muestra es representativa de una poblaci\'on si toda unidad de observaci\'on podr\'a aparecer en la muestra, y esto con una probabilidad conocida.


  \begin{figure}
    \centering
    \includegraphics[width=0.5\linewidth]{../Imagenes/image.png}
\end{figure}   
 \end{frame}


\begin{frame}{¿Por qué muestrear la población?}

    \begin{itemize}
        \item Imposibilidad física de chequear todos los ítems en la población.
        \item El costo de estudiar todos los elementos en la población.
        \item Los resultados del muestreo son normalmente adecuados.
    \end{itemize}
\end{frame}

\begin{frame}{Objetivos del muestreo}

{\color{blue} {\bf Estimar los parámetros de la población.}}

\begin{itemize}
\item {\bf Variable:} lo que se desea medir. Ejemplo: educación, ingreso, empleo.
\item {\bf Parámetro:} valor que describe a la población completa. Ejemplo: promedio poblacional.
\item {\bf Estadístico:} valor calculado con la muestra. Ejemplo: promedio muestral.
\item Como los parámetros suelen ser desconocidos, los estimamos usando estadísticos.
\end{itemize}     
\end{frame}

\begin{frame}{Muestreo aleatorio} 

\begin{itemize}
    \item {\bf Muestreo aleatorio simple:} cada elemento de la población tiene la misma probabilidad de ser incluido.
    \item {\bf Muestreo aleatorio sistemático:} se ordena la población, se elige un punto inicial al azar y luego se selecciona cada $k$-ésimo elemento.
    \item Es simple de aplicar, pero puede fallar si el orden de la lista tiene algún patrón relevante.
\end{itemize}
 
\end{frame}
 

\begin{frame}{Muestreo aleatorio}
\begin{itemize}
    \item {\bf Muestreo aleatorio estratificado:} La población se divide en subgrupos denominados estratos y se selecciona al azar una muestra de cada estrato.
    \item {\bf  Muestreo por conglomerados (cluster):} La población se divide en conglomerados a partir de límites geográficos o de otra clase. Luego se seleccionan los conglomerados al azar y se toma una muestra aleatoria con elementos de cada grupo. 
\end{itemize}
 \end{frame}   



\begin{frame}{Muestreo Estratificado vs. Muestreo por Conglomerados}
\begin{columns}
\column{0.48\textwidth}
\centering
\textbf{Muestreo Estratificado} \\[1ex]
\begin{tikzpicture}
% Estratos homogéneos (una fila por color)
\foreach \i/\color in {0/red, 1/blue, 2/green} {
    \foreach \j in {0,1,2} {
        \filldraw[fill=\color, draw=black] (\j*0.6,\i*0.6) circle (0.2);
    }
    \draw[thick] (-0.3,\i*0.6-0.3) rectangle (1.5,\i*0.6+0.3);
}
\end{tikzpicture}

\vspace{1ex}
\scriptsize
Cada estrato es homogéneo. Se selecciona muestra dentro de cada uno.

\column{0.48\textwidth}
\centering
\textbf{Muestreo por Conglomerados} \\[1ex]
\begin{tikzpicture}
% Conglomerados heterogéneos
\foreach \i/\x in {0/0, 1/2, 2/4} {
    \fill[red] (\x,0) circle (0.2);
    \fill[blue] (\x+0.3,0.3) circle (0.2);
    \fill[green] (\x+0.6,0) circle (0.2);
    \fill[blue] (\x+0.3,-0.3) circle (0.2);
    \draw[black,thick] (\x-0.2,-0.5) rectangle (\x+0.8,0.5);
}
\end{tikzpicture}

\vspace{1ex}
\scriptsize
Cada conglomerado es heterogéneo. Se seleccionan grupos completos.
\end{columns}

\vspace{1ex}
\begin{block}{Comparación Clave}
\begin{itemize}
    \item \textbf{Estratos:} Internamente homogéneos. Se muestrea dentro de cada grupo.
    \item \textbf{Conglomerados:} Internamente heterogéneos. Se seleccionan grupos enteros.
\end{itemize}
\end{block}
\end{frame}

\begin{frame}{Muestreo no probabilístico}
\begin{itemize}
\item La inclusión en la muestra depende del criterio de quien selecciona los casos.
\item Ejemplos: encuestas por conveniencia, voluntarios, entrevistas a personas disponibles.
\item Puede ser útil exploratoriamente, pero limita la inferencia estadística hacia la población.
\end{itemize}
\end{frame}
% \begin{frame}{Muestreo}
%   \begin{itemize}
%       \item  Universo. Es el conjunto de todos los datos u observaciones de un suceso.
% El universo puede ser finito, infinito contable e infinito no numerable.
% \item Poblaci\'on. Elementos del universo que tienen una caracter\'istica com\'un y
% es la que tratamos de estudiar.
% \item Marco Muestral: es el conjunto de unidades del cual se seleccionar\'a una
% muestra.  
% \item Muestra Aleatoria (m.a.). Es un subconjunto de la poblaci\'on donde sus
% elementos son escogidos al azar, es decir, no existe relaci\'on en la elecci\'on
% de los elementos.
%   \end{itemize}

% \end{frame}

% \begin{frame}{Muestreo}
%      Una muestra est\'a correctamente seleccionada cuando es sacada de manera aleatoria
% a partir de toda la poblaci\'on y es suficientemente grande para tener una precisi\'on aceptable.

% \begin{itemize}
%     \item Que no tenga sesgo, las características de la muestra no sobreestimen o no
% subestimen las características de la poblaci\'on que se pretende evaluar.
% \item Todo elemento de la poblaci\'on tiene la posibilidad de ser elegido en la muestra. 
% \end{itemize}
%  \end{frame}

\begin{frame}{Descriptiva vs. inferencial}
   \begin{itemize}
       \item {\color{blue}Estad\'istica Descriptiva o Deductiva.}  
       
       Resume y ordena la información observada en la muestra, sin necesariamente generalizar a toda la población.
\item {\color{blue}Estad\'istica Inductiva o Inferencial.} 

Usa una muestra para inferir conclusiones estadísticamente válidas sobre la población.
   \end{itemize} 
\end{frame}

\begin{frame}{Ideas clave de la clase}
\begin{itemize}
    \item La econometría conecta teoría, datos y estadística.
    \item Correlación no implica causalidad: necesitamos supuestos claros.
    \item El tipo de dato define qué preguntas podemos responder.
    \item Una buena muestra permite aprender sobre una población más grande.
\end{itemize}
\end{frame}





\end{document}
